\chapter{Introducción}
\label{cap:intro}

El \emph{péndulo invertido sobre carro} (carro-péndulo o péndulo lineal) es el banco de pruebas clásico de la teoría de control. Reúne en un sistema pequeño casi todas las dificultades que aparecen en la práctica: un equilibrio inestable, una no linealidad fuerte, un sistema subactuado (dos grados de libertad y una sola entrada), un cero de fase no mínima, saturación del actuador, límites físicos de recorrido y sensores que solo miden posiciones.

Este manual acompaña al paquete de simulación \code{linear\_pendulum} y está organizado en dos partes:

\begin{itemize}
  \item \textbf{Parte teórica} (capítulos~\ref{cap:modelo}--\ref{cap:control}). Se deriva el modelo, se linealiza, se analizan sus propiedades (estabilidad, controlabilidad, observabilidad, fase no mínima), se estudia la instrumentación y se desarrollan cuatro estrategias de control: LQR, PID+PD, \emph{swing-up} por energía y grúa con anti-balanceo.
  \item \textbf{Parte práctica} (capítulos~\ref{cap:simulink}--\ref{cap:practicas}). Se describe el modelo de Simulink, la interfaz de código \code{CartPole} y la interfaz gráfica \code{cp\_gui}. El manual termina con diez prácticas guiadas.
\end{itemize}

Todas las figuras de resultados de este manual se han obtenido simulando el modelo no lineal completo, con actuador, cuantización de los encoders y controlador discreto, y usando los parámetros por defecto de \code{cp\_params.m}. Por tanto, deberían reproducirse con el software tal y como se entrega.

\section{El sistema físico}

La figura~\ref{fig:sistema} muestra el montaje. Un carro de masa $M$ se desplaza sobre un raíl horizontal de longitud útil $2x_{\lim}=\SI{1.6}{m}$. Un actuador lineal (por ejemplo, un motor con correa o un motor lineal) ejerce sobre él una fuerza $F$. En el carro va articulada una varilla uniforme de masa $m$ y longitud $L_p$, que gira libremente en el plano vertical. Dos encoders miden la posición del carro $x$ y el ángulo de la varilla $\theta$.

\begin{figure}[H]
\centering
\begin{tikzpicture}[scale=1.5,>=Latex]
  % raíl
  \fill[pattern=north east lines,pattern color=gray!60] (-4,-0.62) rectangle (4,-0.5);
  \draw[thick] (-4,-0.5) -- (4,-0.5);
  \draw[line width=3pt] (-3.7,-0.5) -- (-3.7,0.1) node[above,font=\small]{tope};
  \draw[line width=3pt] (3.7,-0.5) -- (3.7,0.1) node[above,font=\small]{tope};
  % carro
  \draw[fill=c1!60,rounded corners=2pt] (-0.7,-0.4) rectangle (0.7,0.3);
  \node[font=\small] at (0,-0.1) {$M$};
  \draw[fill=gray] (-0.42,-0.45) circle (0.06); \draw[fill=gray] (0.42,-0.45) circle (0.06);
  % varilla
  \def\ang{25}
  \draw[line width=5pt,c2!80,line cap=round] (0,0.15) -- ++(90-\ang:3.0) coordinate (tip);
  \coordinate (piv) at (0,0.15);
  \path (piv) -- (tip) coordinate[pos=0.5] (cm);
  \fill[black] (cm) circle (0.07);
  \node[right=3pt,font=\small] at (cm) {CdM, $m$, $J$};
  \draw[fill=white] (piv) circle (0.06);
  % ejes y ángulo
  \draw[dashed,gray] (piv) -- ++(0,3.1) node[above,black,font=\small]{vertical};
  \draw[->,thick] (piv) ++(0,1.3) arc (90:90-\ang:1.3);
  \node[font=\small] at ($(piv)+(0.28,1.45)$) {$\theta$};
  % cotas l y Lp
  \draw[<->,gray!80!black] ($(piv)+(-0.25,0)+(180-\ang:0)$) -- ++(90-\ang:1.5) node[midway,left,font=\small]{$l$};
  \draw[<->] ($(piv)+(0.35,-0.15)$) ++(90-\ang:0) -- ++(90-\ang:3.0);
  \node[font=\small,right] at ($(piv)!0.75!(tip)+(0.35,-0.1)$) {$L_p=2l$};
  % fuerzas
  \draw[->,line width=1.6pt,c4] (-2.0,-0.05) -- (-0.75,-0.05) node[midway,above]{$F$};
  \draw[->,line width=1.4pt,c5!70!black] ($(tip)+(-1.1,0)$) -- ($(tip)+(-0.12,0)$) node[pos=0.3,above]{$F_d$};
  \draw[->,thick,gray!70!black] ($(cm)+(0.9,0)$) -- ++(0,-0.7) node[right,font=\small]{$m g$};
  % coordenada x
  \draw[->,thick] (-3.6,-1.0) -- (3.6,-1.0) node[right]{$x$};
  \draw (0,-0.95) -- (0,-1.05) node[below,font=\small]{$x$};
  \draw (-3.7,-0.95) -- (-3.7,-1.05) node[below,font=\small]{$-x_{\lim}$};
  \draw (3.7,-0.95) -- (3.7,-1.05) node[below,font=\small]{$x_{\lim}$};
  % encoders
  \node[draw,circle,fill=yellow!40,inner sep=1pt,font=\scriptsize] at (0.45,0.05) {E$_\theta$};
  \node[draw,rectangle,fill=yellow!40,inner sep=2pt,font=\scriptsize,anchor=north] at (-2.6,-0.62) {Encoder lineal E$_x$};
\end{tikzpicture}
\caption{Esquema del péndulo lineal y convención de signos. $\theta=0$ corresponde a la varilla vertical hacia arriba; $\theta>0$ indica inclinación hacia $+x$.}
\label{fig:sistema}
\end{figure}

\section{Convenciones}

\begin{table}[H]
\centering\small
\begin{tabular}{@{}lll@{}}
\toprule
Símbolo & Significado & Unidad\\
\midrule
$x$ & posición del carro (positiva a la derecha) & m\\
$\theta$ & ángulo de la varilla respecto a la vertical superior; $\theta=\pi$ colgando & rad\\
$F$ & fuerza aplicada por el actuador sobre el carro & N\\
$F_d$ & fuerza de perturbación horizontal en la punta de la varilla & N\\
$\vect{s}$ & estado $[x,\ \xd,\ \theta,\ \thd]^\top$ & --\\
$\vect{y}$ & medida $[x_m,\ \theta_m]^\top$ (solo posiciones) & m, rad\\
$r$ & referencia de posición del carro & m\\
\bottomrule
\end{tabular}
\caption{Convenciones de signos y variables.}
\end{table}

\section{Contenido del paquete de software}

\begin{table}[H]
\centering\small
\begin{tabularx}{\linewidth}{@{}lX@{}}
\toprule
Archivo & Función\\
\midrule
\code{cp\_params.m} & Todos los parámetros físicos, de sensores, temporales y de los controladores.\\
\code{cp\_linearize.m} & Matrices $(A,B,C,D)$ en torno a $\theta_0=0$ o $\theta_0=\pi$.\\
\code{cp\_lqr.m}, \code{cp\_design.m} & Cálculo de las ganancias LQR (con o sin Control System Toolbox).\\
\code{cp\_plant\_block.m} & Dinámica no lineal (código del bloque \emph{MATLAB Function} de la planta).\\
\code{cp\_controller\_block.m} & Controlador discreto con 5 modos (código del bloque del controlador).\\
\code{cp\_ref\_block.m} & Generador de referencias.\\
\code{cp\_build\_model.m} & Genera el modelo \code{cartpole\_sim.slx} mediante un script.\\
\code{CartPole.m} & Clase con la interfaz de código: ensayos por lotes y control en vivo.\\
\code{cp\_gui.m} & Interfaz gráfica en vivo.\\
\code{cp\_draw.m}, \code{cp\_animate.m}, \code{cp\_plot\_results.m} & Animación y gráficas.\\
\code{examples.m} & Recorrido guiado por secciones.\\
\code{python\_reference\_sim.py} & Réplica independiente en Python usada para validar el modelo.\\
\code{validation/} & Scripts de validación cruzada Simulink--Python (apéndice~\ref{cap:validacion}).\\
\bottomrule
\end{tabularx}
\caption{Archivos del paquete \code{linear\_pendulum}.}
\label{tab:archivos}
\end{table}

\section{Requisitos y puesta en marcha}

Se necesita MATLAB y Simulink R2020b o posterior. La Control System Toolbox es opcional: si no está instalada, \code{cp\_lqr} resuelve la ecuación de Riccati por el método del Hamiltoniano (sección~\ref{sec:hamiltoniano}).

\begin{lstlisting}[caption={Primeros pasos.}]
cd src/linear_pendulum
cp_build_model                 % crea cartpole_sim.slx (solo la primera vez)
cp = CartPole();               % interfaz de código
res = cp.run(6);               % ensayo LQR por defecto (theta0 = 0.1 rad)
cp_plot_results(res); cp.animate(res);
cp_gui(cp);                    % interfaz gráfica en vivo
\end{lstlisting}
